\documentclass{article}
\usepackage{amsmath,amssymb,amsthm}

\newtheorem{definition}{Definition}[section]
\newtheorem{theorem}[definition]{Theorem}
\newtheorem{lemma}[definition]{Lemma}
\newtheorem*{remark}{Remark}

\newcommand{\N}{\mathbb{N}}

\title{Divisibility and even numbers}
\author{An example for the autoformalization template}
\date{}

\begin{document}
\maketitle

\section{Divisibility}

We work with the natural numbers $\N$, with their addition and multiplication.

\begin{definition}\label{def:divides}
Let $a$ and $b$ be natural numbers. We say that $a$ \emph{divides} $b$,
written $a \mid b$, if there is a natural number $c$ such that $b = ac$.
\end{definition}

\begin{theorem}\label{thm:refl}
Every natural number divides itself.
\end{theorem}

\begin{proof}
Let $a$ be a natural number. Since $a = a \cdot 1$, the number $c = 1$
shows that $a \mid a$.
\end{proof}

\begin{theorem}\label{thm:trans}
Let $a$, $b$ and $c$ be natural numbers. If $a \mid b$ and $b \mid c$,
then $a \mid c$.
\end{theorem}

\begin{proof}
If $b = ad$ and $c = be$, then $c = (ad)e = a(de)$.
\end{proof}

\section{Even numbers}

\begin{definition}\label{def:even}
A natural number $n$ is \emph{even} if $2 \mid n$.
\end{definition}

\begin{lemma}\label{lem:sum}
If $m$ and $n$ are even natural numbers, then $m + n$ is even.
\end{lemma}

\begin{remark}
The converse fails: $1 + 1$ is even, but $1$ is not.
\end{remark}

\end{document}
